\exercisesection{深度与同调维数}

\begin{enumerate}
\def\labelenumi{\arabic{enumi})}
\tightlist
\item
  设 \(\rho: A \to B\) 为 Noether 局部环之间的一个局部同态，\(N\) 为一个有限生成 B-模。假设 A-模 \(B\) 与 \(N\) 都是平坦的。
\end{enumerate}

\begin{enumerate}
\def\labelenumi{\alph{enumi})}
\item
  证明等式 \(\mathrm{dp}_{\mathrm{B}}(\mathrm{N}) = \mathrm{dp}_{\kappa_{\mathrm{A}} \otimes_{\mathrm{A}} \mathrm{B}}(\kappa_{\mathrm{A}} \otimes_{\mathrm{A}} \mathrm{N})\) （考虑 N 的一个由有限生成自由 B-模组成的分解）。
\item
  设 M 为一个有限生成 A-模。证明等式
\end{enumerate}

\[
\mathrm{dp} _ {\mathrm{B}} (\mathrm{M} \otimes_ {\mathrm{A}} \mathrm{N}) = \mathrm{dp} _ {\mathrm{A}} (\mathrm{M}) + \mathrm{dp} _ {\mathrm{B}} (\mathrm{N}).
\]

（设 K（相应地，L）为 A-模 M（相应地，B-模 N）的由有限生成自由模组成的分解；证明 \(K \otimes_{A} L\) 是 \(M \otimes_{A} N\) 的一个由自由 B-模组成的分解。若 dp \(_{A}\) (M) = m 且 dp \(_{B}\) (N) = n，计算 \(\mathrm{H}_{m+n}(\kappa_{\mathrm{B}} \otimes_{\mathrm{B}} (\mathrm{K} \otimes_{\mathrm{A}} \mathrm{L}))\) ，并用 \(\mathrm{H}_{m}(\kappa_{\mathrm{A}} \otimes_{\mathrm{A}} \mathrm{K})\) 与 \(\mathrm{H}_{n}(\kappa_{\mathrm{B}} \otimes_{\mathrm{B}} \mathrm{L})\) 来表达。

\begin{enumerate}
\def\labelenumi{\arabic{enumi})}
\setcounter{enumi}{1}
\tightlist
\item
  设 A 为 Noether 环，M 为一个 A-模，它具有一个由有限生成自由模组成的有限自由分解 (L, p)。
\end{enumerate}

\begin{enumerate}
\def\labelenumi{\alph{enumi})}
\tightlist
\item
  证明对 A 的每个素理想 p，有
\end{enumerate}

\[
\sum_ {i \geqslant 0} (- 1) ^ {i} \operatorname{rg} _ {\mathrm{A}} (\mathrm{L} _ {i}) = \sum_ {i \geqslant 0} (- 1) ^ {i} \left[ \operatorname{Tor} _ {i} ^ {\mathrm{A}} (\mathrm{M}, \kappa (\mathfrak {p})): \kappa (\mathfrak {p}) \right] = \sum_ {i \geqslant 0} (- 1) ^ {i} \left[ \operatorname{Ext} _ {\mathrm{A}} ^ {i} (\mathrm{M}, \kappa (\mathfrak {p})): \kappa (\mathfrak {p}) \right];
\]

用 \(\chi(\mathrm{M})\) 表示这个整数。

\begin{enumerate}
\def\labelenumi{\alph{enumi})}
\setcounter{enumi}{1}
\item
  证明有 \(\chi(M) \geqslant 0\) （考虑一个与 A 相伴的理想 \(\mathfrak{p}\) ）。
\item
  证明下列条件等价：
\end{enumerate}

\begin{enumerate}
\def\labelenumi{(\roman{enumi})}
\item
  \(\chi (\mathbf{M}) = 0\)
\item
  M 的零化子不化为零；
\item
  M 的零化子含有一个可消去元素。
\end{enumerate}

（若 \(\chi(M) = 0\) ，证明 \(M_{\mathfrak{p}}\) 为零（对每个 \(\mathfrak{p} \in \operatorname{Ass}(A)\) ）；若 \(\chi(M) > 0\) ，类似地证明 Ann(M) 的零化子含有一个可消去元素，这蕴含 Ann(M) = 0。）

\begin{enumerate}
\def\labelenumi{\arabic{enumi})}
\setcounter{enumi}{2}
\tightlist
\item
  设 A 为 Macaulay 局部环。回忆（A, II, 第 186 页，习题 16）：若 A 的一个理想异于 A 且不是严格包含它的两个理想的交，则称它在 A 中是不可约的。证明下列条件等价：
\end{enumerate}

\begin{enumerate}
\def\labelenumi{(\roman{enumi})}
\item
  A 是 Gorenstein 环；
\item
  由一个极大割序列生成的 A 的每个理想都是不可约的；
\item
  在 A 中存在一个生成不可约理想的极大割序列。
\end{enumerate}

（化为 A 的维数为 0 的情形，并考虑 A 的底座 \(\mathrm{Hom}_{\mathrm{A}}(\kappa_{\mathrm{A}},\mathrm{A})\) 。）

\begin{enumerate}
\def\labelenumi{\arabic{enumi})}
\setcounter{enumi}{3}
\tightlist
\item
  设 A 为 Noether 环，M 为一个非零的有限生成 A-模。有 grade(M) ≤ dpA(M)；若等号成立，则称 M 是完满的。
\end{enumerate}

\begin{enumerate}
\def\labelenumi{\alph{enumi})}
\tightlist
\item
  设 d 为整数；证明下列条件等价：
\end{enumerate}

\begin{enumerate}
\def\labelenumi{(\roman{enumi})}
\item
  A-模 M 是投射维数为 d 的完满模；
\item
  有 \(\operatorname{Ext}_{\mathrm{A}}^{p}(\mathrm{M},\mathrm{A}) = 0\) （对 \(p\neq d\) ）
\item
  对 M 的支撑集中的每个素理想（相应地，极大理想）\(\mathfrak{p}\) ，\(\mathrm{A}_{\mathfrak{p}}\) -模 \(\mathrm{M}_{\mathfrak{p}}\) 是投射维数为 \(d\) 的完满模。
\end{enumerate}

（先证明 (i) 与 (ii) 的等价性。）

\begin{enumerate}
\def\labelenumi{\alph{enumi})}
\setcounter{enumi}{1}
\item
  设 J 为 A 的由一个完全割序列生成的理想，并设 m 为一个整数 \(\geqslant 1\) 。证明 A-模 \(A/J^{m}\) 是完满的。
\item
  若 M 是完满的，则它的相伴素理想是使 \(\text{prof}(A_{\mathfrak{p}}) = \text{dp}_{\mathrm{A}}(M)\) 的那些理想 p。
\item
  对每个投射维数为 d 的完满模 N，令 \(\mathrm{N}^{\prime} = \mathrm{Ext}_{\mathrm{A}}^{d}(\mathrm{N}, \mathrm{A})\) 。证明 \(N^{\prime}\) 是完满的，其投射维数为 d，且 \((\mathrm{N}^{\prime})^{\prime}\) 同构于 N（考虑 N 的一个投射分解）。我们有 \(\mathrm{Ass}(\mathrm{N}^{\prime}) = \mathrm{Ass}(\mathrm{N})\) 。
\item
  设 \(x\) 为 A 的根中的一个非零因子，使得乘法映射 \(x_{\mathrm{M}}\) 是单射。为使 A-模 M 是投射维数为 \(d\) 的完满模，必要且充分的是 (A/xA)-模 M/xM 也是完满的。
\item
  投射维数有限的 Macaulay 模是完满的（化为局部情形，并对 dim(M) 用归纳法论证）。
\item
  若 A 是 Macaulay 环，则完满模恰是投射维数有限的那些 Macaulay 模。
\end{enumerate}

\begin{enumerate}
\def\labelenumi{\arabic{enumi})}
\setcounter{enumi}{4}
\tightlist
\item
  \begin{enumerate}
  \def\labelenumii{\alph{enumii})}
  \tightlist
  \item
    设 \(k\) 为一个域，R 为一个有限次的 \(k\) -代数。为使 R 是 Gorenstein 环，必要且充分的是 R-模 \(\mathrm{Hom}_{k}(\mathrm{R}, k)\) 是秩为 1 的自由模（化为 R 是局部的情形，并注意 R-模 \(\mathrm{Hom}_{k}(\mathrm{R}, k)\) 是内射的且含有一个同构于 \(\kappa_{\mathrm{R}}\) 的子模；另一个证明见 § 8, 第 6 节，命题 5 的推论）。
  \end{enumerate}
\end{enumerate}

\begin{enumerate}
\def\labelenumi{\alph{enumi})}
\setcounter{enumi}{1}
\item
  假设上述条件成立，且代数 R 是 N 型分次的，其中 \(R_{0} = k\) 。设 s 为使 \(R_{s} \neq 0\) 的最大整数。证明 \(R_{s}\) 在 k 上的维数为 1，且等于 R 的底座；对每个非零的 k-线性形式 \(\ell\) （在 \(R_{s}\) 上），k-双线性形式 \((x, y) \mapsto \ell(xy)\) （在 \(R_{i} \times R_{s- i}\) 上）对每个 i 诱导从 \(R_{i}\) 到 \(\operatorname{Hom}_{k}(R_{s-i}, k)\) 的一个同构。
\item
  设 A 为 Noether 环，B 为一个作为有限型平坦模的 A-代数。证明下列条件等价：
\end{enumerate}

\begin{enumerate}
\def\labelenumi{(\roman{enumi})}
\item
  B-模 \(\operatorname{Hom}_{\mathrm{A}}(\mathbf{B},\mathbf{A})\) 是秩为 1 的投射模；
\item
  对 A 的每个极大理想（相应地，素理想）\(\mathfrak{p}\) ，环 \(\kappa(\mathfrak{p})\otimes_{\mathrm{A}}\mathrm{B}\) 是 Gorenstein 环。
\end{enumerate}

当 B 是 Gorenstein 环时，这些条件特别地被满足。

¶ 6) 设 A 为 Noether 局部环，M 为一个具有有限投射维数的有限生成 A-模，N 为一个 A-模，使得 \(\operatorname{Tor}_{i}^{\mathrm{A}}(\mathrm{M},\mathrm{N}) = 0\) （对 i \textgreater{} 0）。

\begin{enumerate}
\def\labelenumi{\alph{enumi})}
\tightlist
\item
  对每个整数 n 定义一个典范同构
\end{enumerate}

\[
\operatorname{Ext} _ {\mathrm{A}} ^ {n} (\kappa_ {\mathrm{A}}, \mathrm{M} \otimes_ {\mathrm{A}} \mathrm{N}) \longrightarrow \bigoplus_ {p \geqslant n} \operatorname{Tor} _ {p - n} ^ {\mathrm{A}} (\kappa_ {\mathrm{A}}, \mathrm{M}) \otimes_ {k} \operatorname{Ext} _ {\mathrm{A}} ^ {p} (\kappa_ {\mathrm{A}}, \mathrm{N}),
\]

并证明有 \(\mathrm{di}_{\mathrm{A}}(\mathrm{M}\otimes_{\mathrm{A}}\mathrm{N})\leqslant \mathrm{di}_{\mathrm{A}}(\mathrm{N})\)。

（设 (P,p) 为 M 的一个有限自由分解，(I,e) 为 N 的一个内射分解，并设 C 为复形 \(P \otimes_A I\) 。证明复形

\[
0 \to \mathrm{C} ^ {0} / \mathrm{B} ^ {0} (\mathrm{C}) \to \mathrm{C} ^ {1} \to \mathrm{C} ^ {2} \to \dots
\]

定义 \(M \otimes_A N\) 的一个内射分解。）

\begin{enumerate}
\def\labelenumi{\alph{enumi})}
\setcounter{enumi}{1}
\item
  证明不等式 \(\mathrm{prof}_{\mathrm{A}}(\mathrm{N}) = \mathrm{dp}_{\mathrm{A}}(\mathrm{M}) + \mathrm{prof}(\mathrm{M} \otimes_{\mathrm{A}} \mathrm{N})\) 。由此导出 Auslander-Buchsbaum 定理的另一个证明。
\item
  进而假设 N 具有有限投射维数。证明等式 \(\mathrm{dp}_{\mathrm{A}}(\mathrm{M}\otimes_{\mathrm{A}}\mathrm{N})=\mathrm{dp}_{\mathrm{A}}(\mathrm{M})+\mathrm{dp}_{\mathrm{A}}(\mathrm{N})\) 。
\item
  假设 A 是 Gorenstein 环，并用 d 表示它的维数。设 n 为整数；证明 \(\kappa_{A}\) -向量空间 \(\operatorname{Ext}_{\mathrm{A}}^{n}(\mathbf{k}_{\mathrm{A}},\mathrm{M})\) 、\(\operatorname{Tor}_{d-n}^{\mathrm{A}}(\mathbf{k}_{\mathrm{A}},\mathrm{M})\) 与 \(\operatorname{Ext}_{\mathrm{A}}^{d-n}(\mathrm{M},\mathbf{k}_{\mathrm{A}})\) 具有相同的维数。
\end{enumerate}

\begin{enumerate}
\def\labelenumi{\arabic{enumi})}
\setcounter{enumi}{6}
\tightlist
\item
  设 A 为 Noether 局部环，M 为一个有限生成 A-模，T 为一个具有有限内射维数的有限生成 A-模。证明等式
\end{enumerate}

\[
\operatorname{prof} (\mathrm{A}) = \operatorname{prof} _ {\mathrm{A}} (\mathrm{M}) + \sup \left\{i \mid \operatorname{Ext} _ {\mathrm{A}} ^ {i} (\mathrm{M}, \mathrm{T}) \neq 0 \right\}
\]

（对 \(\mathrm{prof}_{\mathrm{A}}(\mathrm{M})\) 用归纳法推理，并利用命题 9）。

\begin{enumerate}
\def\labelenumi{\arabic{enumi})}
\setcounter{enumi}{7}
\tightlist
\item
  设 A 为环，并设 \(u: \mathrm{A}^p \to \mathrm{A}^q\) 为一个 A-线性映射。我们称 \(u\) 的秩——并用 \(\text{rg}(u)\) 表示——为最大的整数 \(r\) ，它使 \(\wedge^r u \neq 0\) 。用 \(\mathfrak{d}(u)\) 表示由 \(\text{rg}(u)\) 阶子式（属于 \(u\) 的矩阵）所生成的 A 的理想。
\end{enumerate}

\begin{enumerate}
\def\labelenumi{\alph{enumi})}
\item
  证明为使 \(\mathfrak{d}(u)\) 等于 A，必要且充分的是 \(u\) 的余核是秩为 \(q - \mathrm{rg}(u)\) 的投射模（化为 \(\mathrm{A}\) 是局部的情形，并在一个适当的基中写出 \(u\) 的矩阵）。
\item
  设 \(v: \mathbf{A}^q \to \mathrm{A}^s\) 为一个 A-线性映射；我们假设有 \(\mathfrak{d}(u) = \mathfrak{d}(v) = \mathbf{A}\) 为使 \(\operatorname{Ker} v = \operatorname{Im} u\) ，必要且充分的是有 \(\operatorname{rg}(u) + \operatorname{rg}(v) = q\) （同样的方法）。
\item
  若 \(\operatorname{rg}(u) = q\) ，证明 \(\mathfrak{d}(u)\) 零化 \(u\) 的余核（通过化为 \(p = q\) 的情形）。
\end{enumerate}

¶ 9) 设 A 为环，并设 (L, d) 为一个由有限生成自由 A-模组成的有界复形，它在次数 \textless{} 0 处为零。我们打算证明下列条件等价：

\begin{enumerate}
\def\labelenumi{(\roman{enumi})}
\item
  有 \(\mathrm{H}_i(\mathrm{L}) = 0\) （对 \(i > 0\) ）；
\item
  有 \(\operatorname{prof}_{\mathrm{A}}(\mathfrak{d}(d_i); \mathrm{A}) \geqslant i\) 且 \(\operatorname{rg}(d_i) + \operatorname{rg}(d_{i+1}) = \operatorname{rg}_{\mathrm{A}}(\mathrm{L}_i)\) （对每个 \(i \geqslant 1\) ）。
\end{enumerate}

\begin{enumerate}
\def\labelenumi{\alph{enumi})}
\item
  证明只需对环 A 是 Noether 的情形证明这一命题（考虑由矩阵 \(d_{i}\) （对 \(i \geqslant 1\) ）的系数所生成的 A 的子环）。以下假设这个条件成立。
\item
  若 A 是一个维数为 0 的局部环，证明 (ii) 蕴含 (i)（利用习题 8, b)）。
\item
  假设环 A 是局部的，条件 (ii) 成立，且 \(\operatorname{Supp}(\mathrm{H}_i(\mathrm{L})) = \{\mathfrak{m}_{\mathrm{A}}\}\) （对 \(i \geqslant 1\) ）。证明 \(\mathrm{H}_i(\mathrm{L}) = 0\) （对 \(i \geqslant 0\) ）（设 \(p\) 为使 \(\mathfrak{d}(d_p) \neq \mathrm{A}\) 的最大整数；由习题 8, b) 推出 \(\mathrm{H}_i(\mathrm{L}) = 0\) （对 \(i > p\) ），然后由习题 8, a) 以及命题 3 的推论 2 （ \(n^{\circ} 2\) ）推出复形 \(0 \to L_p / B_p(L) \to L_{p-1} \to \ldots \to L_0\) 在次数 \(>0\) 处是零调的）。
\item
  证明蕴含关系 (ii)⇒(i)（化为 A 是局部的情形，并对 dim(A) 用归纳法论证）。
\item
  我们假设复形 L 在次数 \textgreater{} 0 处是零调的。设 p 为与 A 相伴的一个素理想。证明对每个 \(i \geqslant 1\) ，(Coker \(d_{i}\) )p 是一个自由 Ap-模，其秩为 \(\sum_{p \geqslant i-1} (-1)^{p} \operatorname{rg}_{\mathrm{A}}(\mathrm{L}_{p})\) ；此外，我们有 \(\operatorname{rg}((d_{i})_{\mathfrak{p}}) = \sum_{p \geqslant i} (-1)^{p-i} \operatorname{rg}_{\mathrm{A}}(\mathrm{L}_{p})\) 以及 \(\mathfrak{d}((d_{i})_{\mathfrak{p}}) = A_{\mathfrak{p}}\) （习题 8）。
\end{enumerate}

\(f)\) 于是我们有 \(\operatorname{rg}(d_i) = \sum_{p \geqslant i} (-1)^{p - i} \operatorname{rg}_A(L_p)\) ，从而有 \(\operatorname{rg}(d_i) + \operatorname{rg}(d_{i+1}) = \operatorname{rg}_\mathrm{A}(L_i)\) （对每个 \(i \geqslant 1\) ）。

\(g)\) 设 \(i\) 为一个整数 \(\geqslant 1\) ，并设 \(\mathfrak{p}\) 为 \(\mathrm{A}\) 的一个素理想。证明条件 \(\mathrm{dp}_{\mathrm{A}_{\mathfrak{p}}}(H_0(L)_{\mathfrak{p}}) \geqslant i\) 等价于 \(\mathfrak{d}(d_i)A_{\mathfrak{p}} \neq A_{\mathfrak{p}}\) ，亦即 \(\mathfrak{d}(d_i) \subset \mathfrak{p}\) （利用习题 8, a)）。利用 § 1, 第 5 节的命题 8 以及定理 1，推出有 \(\operatorname{prof}_{\mathrm{A}}(\mathfrak{d}(d_i); A) \geqslant i\) ，这就证明了蕴含关系 (i) \(\Rightarrow\) (ii)。

¶ 10) 我们保留习题 9 的假设；我们假设复形 L 在次数 ≠ 0 处是正合的。

\begin{enumerate}
\def\labelenumi{\alph{enumi})}
\item
  证明包含关系 \(V(\mathfrak{d}(d_{i+1})) \subset V(\mathfrak{d}(d_i))\) （对每个 \(i \geqslant 1\) ）（注意由习题 8，素理想 p 不含 \(\mathfrak{d}(d_i)\) 当且仅当 \(A_{\mathfrak{p}}\)-模 \((\operatorname{Coker} d_i)_{\mathfrak{p}}\) 是自由的）。
\item
  我们现在假设 \(\mathrm{H}_{0}(\mathrm{L})\) 的零化子不约化为零。证明有 \(\mathrm{rg}(d_{1}) = \mathrm{rg}_{\mathrm{A}}(\mathrm{L}_{0})\) （若 f 是 \(\mathfrak{o}(d_{1})\) 的一个可消去元素，则 \(A_{f}\) -模 \(\mathrm{H}_{0}(\mathrm{L})_{f}\) 依习题 8 是具有常秩的投射模，故必为零）；由此推出 \(\mathrm{H}_{0}(\mathrm{L})\) 的支撑集是 \(\mathrm{V}(\mathfrak{o}(d_{1}))\) 。
\item
  我们置 \(p = \text{prof}_A(\mathfrak{d}(d_1); A)\) 。证明有 \(V(\mathfrak{d}(d_i)) = \text{Supp}(\mathbf{H}_0(L))\) （对 \(1 \leqslant i \leqslant p\) ）（如 \(a\) ) 的证明中那样推理，并注意：若一个素理想 \(\mathfrak{p}\) （属于 \(\mathrm{H}_0(L)\) 的支撑集）不含 \(\mathfrak{d}(d_k)\) ，则将有 \(\text{dp}_{\mathrm{A}_{\mathfrak{p}}} \, H_0(L)_{\mathfrak{p}} < k \leqslant \text{grade}_{\mathrm{A}_{\mathfrak{p}}} \, H_0(L)_{\mathfrak{p}}\) 。
\end{enumerate}

\begin{enumerate}
\def\labelenumi{\arabic{enumi})}
\setcounter{enumi}{10}
\tightlist
\item
  设 A 为局部环。
\end{enumerate}

\begin{enumerate}
\def\labelenumi{\alph{enumi})}
\tightlist
\item
  设 n 为一个整数 \(\geqslant 1\) ，\(u : A^{n-1} \to A^n\) 为一个 A-线性映射，\(\mathfrak{d}(u)\) 为由 u 的 n-1 阶子式所生成的理想（习题 8）。若 \(\operatorname{prof}_{\mathrm{A}}(\mathfrak{d}(u); A) \geqslant 2\) ，证明 \(\mathfrak{d}(u)\) 容许分解
\end{enumerate}

\[
0 \rightarrow \mathrm{A} ^ {n - 1} \xrightarrow {u} \mathrm{A} ^ {n} \xrightarrow {v} \mathfrak {d} (u) \rightarrow 0,
\]

其中 \(v\) 由同态 \(\wedge^{n-1}(^t u): A^n \to A\) 诱导（应用习题 9）。

\begin{enumerate}
\def\labelenumi{\alph{enumi})}
\setcounter{enumi}{1}
\tightlist
\item
  反之，设 J 为 A 的一个投射维数为 1 的理想。证明存在 A 的一个可消去元素 a、一个整数 n 以及一个 A-线性映射 \(u: A^{n-1} \to A^n\) ，使得 \(\mathrm{J} = a\mathfrak{d}(u)\) ；有 \(\operatorname{prof}_{\mathrm{A}}(\mathrm{J}; \mathrm{A}) = 2\) （“Hilbert-Burch 定理”：由习题 9 推出 J 的一个极小分解形如 \(0 \to A^{n-1} \xrightarrow{u} \mathrm{A}^n \to J\) ，其中 \(\operatorname{prof}_{\mathrm{A}}(\mathfrak{d}(u); \mathrm{A}) \geqslant 2\) ，然后应用 a) 以及 § 1 的习题 1，并注意 \(\mathfrak{d}(u)\) 含有 A 的一个可消去元素）。
\end{enumerate}

\begin{enumerate}
\def\labelenumi{\arabic{enumi})}
\setcounter{enumi}{11}
\tightlist
\item
  设 A 为 Noether 局部环，并设 M 与 N 为非零的有限生成 \(\mathrm{A}\)-模。假设 M 的投射维数有限。设 \(q\) 为使 \(\operatorname{Tor}_q^A(M, N) \neq 0\) 的最大整数。
\end{enumerate}

\begin{enumerate}
\def\labelenumi{\alph{enumi})}
\item
  假设 \(\mathrm{prof}_{\mathrm{A}}(\mathrm{Tor}_{q}^{\mathrm{A}}(\mathrm{M},\mathrm{N})) = 0\) 。证明等式 \(\mathrm{prof}_{\mathrm{A}}(\mathrm{N}) + q = \mathrm{dp}_{\mathrm{A}}(\mathrm{M})\) （可对 \(\mathrm{prof}_{\mathrm{A}}(\mathrm{N})\) 用归纳法推理，考虑一个元素 \(x\) （属于 \(\mathfrak{m}_{\mathrm{A}}\) ，且在 N 中为非零因子））。
\item
  我们进一步假设 A-模 \(M \otimes_{A} N\) 具有有限长度。由 a) 推出：\(\operatorname{Tor}_{i}^{A}(M, N)\) 对 i \textgreater{} depth(A) - depth(M) - depth(N) 为零，且在等号成立时非零。特别地，\(\operatorname{prof}(M) + \operatorname{prof}(N) \leqslant \operatorname{prof}(A)\) 。
\end{enumerate}

¶ 13) 设 A 为 Noether 局部环。假设下列条件被满足 \(^{1}\) ：

（GM） 对每个 Noether 局部 A-代数 B，存在一个 B-模 M，使得 \(\mathfrak{m}_{\mathrm{B}}\mathrm{M}\neq\mathrm{M}\) 且 \(\operatorname{prof}_{\mathrm{B}}(\mathrm{M})=\dim(\mathrm{B})\) 。

\begin{enumerate}
\def\labelenumi{\alph{enumi})}
\item
  设 P 为一个由平坦 A-模组成的、有限长度为 \(\ell\) 的复形，使得 \(\operatorname{Supp}(\mathrm{H}(\mathrm{P})) = \{\mathfrak{m}_{\mathrm{A}}\}\) 。证明有 \(\ell \geqslant \dim(\mathrm{A})\) （利用 § 1 的习题 6） \(^2\) 。
\item
  设 B 为 Noether 环，\(\rho: A \to B\) 为一个同态，\(^a\rho: \operatorname{Spec}(B) \to \operatorname{Spec}(A)\) 为相应的映射，M 为一个有限生成 A-模。证明 \(^a\rho^{-1}(\operatorname{Supp}(M))\) 在 \(\operatorname{Spec}(B)\) 中的余维数 \(\leqslant \mathrm{dp}_A(M)\) （把 \(a\) ) 应用于复形 \(L \otimes_A B_q\) ，其中 \(q\) 是 \(^a\rho^{-1}(\operatorname{Supp}(M))\) 的一个分支的极小素理想，而 L 是 M 的一个自由分解）。
\item
  设 M、N 为有限生成 A-模，其中 \(M \neq 0\) 。证明不等式
\end{enumerate}

\[
\dim (N) \leqslant \mathrm{dp} _ {A} (M) + \dim (M \otimes_ {A} N)
\]

（相交定理：对整数 \(d = \dim(M \otimes_{\mathrm{A}} N)\) 用归纳法推理。若 d = 0，把 b) 应用于环 B = A/Ann(M)；一般情形考虑 A 的一个对 N 以及对 M \(\otimes_{A} N\) 正则的元素）。

\begin{enumerate}
\def\labelenumi{\alph{enumi})}
\setcounter{enumi}{3}
\tightlist
\item
  为使存在一个具有有限长度且有限投射维数的 A-模，必要且充分的是 A 为 Macaulay 环。
\end{enumerate}

\begin{enumerate}
\def\labelenumi{\arabic{enumi})}
\setcounter{enumi}{13}
\tightlist
\item
  设 \(\mathrm{A}\) 为满足条件 (GM) 的 Noether 局部环（习题 13），M 为一个具有有限投射维数的有限生成 A-模。
\end{enumerate}

\begin{enumerate}
\def\labelenumi{\alph{enumi})}
\tightlist
\item
  证明不等式
\end{enumerate}

\[
\dim (A) \leqslant \mathrm{dp} _ {A} (M) + \dim (M) \quad \dim (A) - \operatorname{prof} (A) \leqslant \dim (M) - \operatorname{prof} (M).
\]

若此外模 M 是完满的（习题 4），则有 \(\dim(A) = \mathrm{dp}_{A}(M) + \dim(M)\) （利用 § 1 的习题 15）。

\begin{enumerate}
\def\labelenumi{\alph{enumi})}
\setcounter{enumi}{1}
\item
  设 \(\mathfrak{p} \in \mathrm{Supp}(M)\) ；若 \(\mathrm{A}_{\mathfrak{p}}\) -模 \(\mathrm{M}_{\mathfrak{p}}\) 是 Macaulay 的，则环 \(\mathrm{A}_{\mathfrak{p}}\) 是 Macaulay 环。特别地，当 \(\mathfrak{p}\) 是 \(\mathrm{Supp}(M)\) 的一个极小素理想时即属此情形。
\item
  设 \(\mathbf{x} = (x_{1},\ldots ,x_{r})\) 为 \(\mathfrak{m}_{\mathrm{A}}\) 中元素的一个正则序列，它生成一个具有有限投射维数的理想 J。证明序列 \(\mathbf{x}\) 对 A 是完全正则的（设 \(\mathfrak{p}\in V(J)\) ，并设 \(\mathfrak{q}\subset \mathfrak{p}\) 为 \(V(J)\) 的一个极小素理想；注意有 \(\mathrm{prof}(A_{\mathfrak{p}})\geqslant \mathrm{dp}_{A_{\mathfrak{p}}}(A_{\mathfrak{p}} / J_{\mathfrak{p}})\geqslant \mathrm{dp}_{A_{\mathfrak{q}}}(A_{\mathfrak{q}} / J_{\mathfrak{q}})\) ，且有 \(\mathrm{dp}_{A_{\mathfrak{q}}}(A_{\mathfrak{q}} / J_{\mathfrak{q}})\geqslant r\) （依 \(b\) ）），从而最终有 \(\mathrm{prof}_{\mathrm{A}}(\mathrm{J};\mathrm{A})\geqslant r\) ；借助 § 1, 第 3 节的定理 1 的推论 1，断定 \(\mathbf{x}\) 是完全正则的）。
\item
  证明每个 M-正则序列都是 A-正则的。（证明 A 的每个零化子非零的元素 a 都不是 M-正则的；为此，通过局部化化为 Supp(M) ∩ Supp(a) = \(\{\mathfrak{m}_{A}\}\) 的情形；把习题 13 a) 应用于复形 \(L \otimes_A A/p\)（其中 L 是 M 的一个自由分解，而 p 是 a 的一个相伴素理想），证明不等式 \(\operatorname{prof}(A) \leqslant \dim(A/p) \leqslant \operatorname{dp}_{A}(M)\) ，并由此推出 prof(M) = 0。）
\item
  为使 A 是整环，必要且充分的是它含有一个投射维数 \(< +\infty\) 的素理想（若 \(\mathrm{dp}_{\mathrm{A}}(\mathfrak{p}) < +\infty\) ，证明 A 可等同于正则环 \(\mathrm{A}_{\mathfrak{p}}\) 的一个子环）。
\end{enumerate}

\begin{enumerate}
\def\labelenumi{\arabic{enumi})}
\setcounter{enumi}{14}
\tightlist
\item
  设 A 为 Noether 局部环，M 为一个具有有限内射维数的有限生成 A-模。
\end{enumerate}

\begin{enumerate}
\def\labelenumi{\alph{enumi})}
\item
  设 \(\mathfrak{p} \in \operatorname{Supp}(M)\) 。证明等式 \(\operatorname{prof}(A) = \operatorname{prof}(A_{\mathfrak{p}}) + \dim(A/p)\) 。（设 \(p = \operatorname{prof}(A_{\mathfrak{p}})\) 且 \(q = \dim(A/p)\) ；注意有 \(p = \operatorname{di}_{A_{\mathfrak{p}}} (M_{\mathfrak{p}})\) ，从而 \(\operatorname{Ext}_{A_{\mathfrak{p}}}^{p}(\kappa(\mathfrak{p}), M_{\mathfrak{p}}) \neq 0\) ，故依 § 1, 第 7 节的引理 3 有 \(\operatorname{Ext}_{A}^{p+q}(\kappa_A, M) \neq 0\) ，这蕴含 \(\operatorname{prof}(A) = \operatorname{di}_A(M) = p + q\) 。）若 \(\operatorname{Supp}(M) = \operatorname{Spec}(A)\) ，则 A 是 Macaulay 环。
\item
  证明等式 grade(M) + dim(M) = prof(A)（利用 a) 以及 § 1 的习题 15）。
\end{enumerate}

¶ 16) 设 A 为 Macaulay 环，M 为一个具有有限投射维数的 A-模。证明为使 M 是平坦的，必要且充分的是：对 A 的每个由完全割序列生成的理想 J，典范同态 \(J \otimes_{A} M \to JM\) 是双射（用对序列长度的归纳法证明这个条件蕴含 \(\operatorname{Tor}_{i}^{A}(A/J, M) = 0\) （对一切 i \textgreater{} 0），然后通过对 n 的递降归纳法证明 \(\operatorname{Tor}_{n}^{A}(N, M) = 0\) （对每个有限生成 A-模 N 及每个整数 n \textgreater{} 0））。例子：Dedekind 环上的无挠模是平坦的（参看 VII, § 2, 习题 14）。
% semantic-fragments: legacy-input-seam
\relax{}
